\section{Distinguish a full queue from an empty one}
Let \(H\) be head, \(T\) tail, \(U\) occupied bytes, and \(C\) capacity. A valid
state satisfies
\[
  0\leq H<C,\qquad0\leq T<C,\qquad0\leq U\leq C,
  \qquad T=(H+U)\bmod C.
\]
Starting at head, the encoded records occupy exactly \(U\) consecutive circular
bytes. Their prefix lengths delimit the records. The caller guarantees that each
record fits the capacity and that the occupied sequence ends at tail.

When \(U=0\), no record is queued. When \(U=C\), every byte is occupied. In both
cases head equals tail. Comparing the two indices alone therefore cannot select
empty versus full under this representation. Some ring designs reserve one slot to
make the distinction another way. This contract instead retains an occupied-byte
count and allows all \(C\) bytes to be used.

For a push of encoded size \(N\), free space is \(C-U\). Success requires
\(N\leq C-U\). The new state is
\[
  H'=H,\qquad T'=(T+N)\bmod C,\qquad U'=U+N.
\]
The free region begins at tail and continues circularly to head. It can appear as
two physical spans while remaining one logical region. A record need not fit wholly
between tail and the physical array end.

For a successful pop of the oldest record, its prefix gives \(L\) and hence
\(N=L+1\). The new state is
\[
  H'=(H+N)\bmod C,\qquad T'=T,\qquad U'=U-N.
\]
Pop changes which bytes are occupied. It does not erase them. Old bytes in free
space have no current record meaning; later pushes may overwrite them. Treating a
stale prefix as a queued record when \(U=0\) would manufacture input that does not
exist.
